\chapter{On-Device Artifact Detection}

\section{Overview}

This chapter presents the on-device detector that powers \TimeLens{}: a compact
YOLOv8n model that recognizes 51 \GEM{} artifacts in real time on a phone. The central
finding is that, for this fine-grained domain, \textbf{label quality dominates model
architecture and input resolution}. We reach a deployable detector not by enlarging
the network but by progressively improving the training labels, from foundation-model
auto-annotation through spatial cleaning rules to full hand annotation.

\section{Problem and Constraints}

Pointing a phone at an exhibit must (i) detect that an artifact is present, (ii)
classify it as one of 51 catalogued artifacts, and (iii) track it with a bounding box
in real time. Deployment constraints rule out cloud detectors: inference must run
on-device (no per-frame round-trip), in real time ($\geq 2.5$\,FPS on mid-tier
phones), with a compact model ($<8$\,MB TFLite asset), and be robust to handheld
capture. These constraints led to YOLOv8-nano with FP16 TensorFlow Lite
deployment~\citep{yolov8,ipalakova2025comparative}.

\section{Objectives and Contributions}

\begin{itemize}[leftmargin=1.5em]
  \item A foundation-model auto-annotation pipeline (YOLO-World) plus four spatial
        sanity rules that recover most labeling failures without manual work.
  \item A controlled iteration study (v1$\rightarrow$v3) that isolates label quality
        from architecture and resolution.
  \item An honest analysis of validation frame leakage and training-to-deployment
        domain shift.
  \item A deployable 5.97\,MB TFLite detector with end-to-end NMS.
\end{itemize}

\section{Methodology}

\subsection{Dataset and Acquisition}
The source material is roughly 50 videos of GEM artifacts captured before public
opening. Frames extracted at about one per second yield approximately 30{,}000
candidate images across 51 classes. Manual annotation of 30k images is infeasible for
a student team, motivating an auto-annotation-first pipeline.

\subsection{v1 --- Auto-Annotated Baseline}
The first labels were produced by YOLO-World (\texttt{yolov8s-worldv2})\citep{yoloworld},
a text-prompted open-vocabulary detector. Each class received a prompt (e.g., ``statue
of Ramesses II'', ``stone stela''), and the highest-confidence box per frame was kept
as ground truth. A YOLOv8n model trained on these labels at $320\times320$ reached
mAP@0.5 $=0.751$, but per-class analysis revealed \textbf{ten classes below mAP@0.5
$=0.30$}. Inspection showed two failure modes: \emph{Mode A} (nine classes), where
YOLO-World boxed the wrong object in the frame (a wall card, a bench, a neighboring
statue); and \emph{Mode B} (one class, \emph{Head of Ramesses III}), sibling-class
confusion between a bust and full-body statues. The architecture was capable; the
labels were the bottleneck.

\subsection{v2 --- Spatial Sanity Rules}
A second-pass annotator added four geometric rules to the YOLO-World output:
(1)~restrict the ``stela'' prompt to the two true stela classes; (2)~drop boxes whose
center lies in the bottom quarter of the frame ($c_y>0.75$), which are almost always
floor signs or bystanders; (3)~pick the largest-area box rather than the
highest-confidence one (artifacts dominate the frame; spurious text boxes are small);
and (4)~for each video, drop boxes whose center deviates from the per-video median by
more than $0.20$. These rules dropped about $6.9\%$ of frames, leaving $27{,}134$.
Retraining (still $320\times320$) raised mAP@0.5 to $0.867$ ($+0.116$ from a labeling
change alone) and reduced broken classes from ten to three. Three classes remained at
zero because their auto-labels were unrecoverable.

\subsection{v2 at 640 and End-to-End NMS}
Raising the input to $640\times640$ and exporting with end-to-end NMS
(\texttt{nms=True}, output \texttt{[1,300,6]}) left mAP@0.5 unchanged at $0.867$ but
lifted mAP@0.5:0.95 from $0.586$ to $0.777$ ($+0.191$): higher resolution improves
\emph{localization} but cannot fix classes whose labels are wrong. The data-quality
ceiling was unmistakable.

\subsection{v3 --- Full Hand Annotation}
All 51 classes were re-annotated by hand in Roboflow ($24{,}384$ images,
$70/20/10$ split), drawing full-artifact boxes, avoiding peripheral objects, and
carefully distinguishing similar artifacts. Warm-started from the v2@640 weights and
trained for 100 epochs on a GTX~1650~Ti, v3 reached precision $0.998$, recall $1.000$,
mAP@0.5 $0.995$, and mAP@0.5:0.95 $0.924$, with \textbf{zero broken classes}. The
validation classification loss collapsed from $0.892$ to $0.166$ --- the model's
quantitative way of saying the labels are now internally consistent.

\section{Core Algorithms}

Two procedures capture the steps most prone to error. The first formalizes the
auto-annotation pass and its spatial sanity rules (the v1$\rightarrow$v2 transition);
the second shows how the exported model's end-to-end NMS output is parsed and
deduplicated on-device.

\begin{algorithm}[htbp]
\caption{Auto-annotation with spatial sanity rules (v1$\rightarrow$v2)}
\KwIn{video frames $F$; per-class text prompts $P$; true-stela class set $S$}
\KwOut{one labeled box per usable frame}
\ForEach{frame $f \in F$ (class $c$ from its folder)}{
  $B \leftarrow \textsc{YoloWorld}(f, P)$\;
  \If{$c \notin S$}{remove stela-prompt detections from $B$}
  remove boxes with center $c_y > 0.75$\;
  \If{$B \neq \emptyset$}{
    $b^\star \leftarrow \arg\max_{b\in B}\;\mathrm{area}(b)$ \tcp*{tie-break on confidence}
    drop $b^\star$ if its center deviates $>0.20$ from the per-video median center\;
    \If{$b^\star$ kept}{emit $(b^\star, c)$}
  }
}
\end{algorithm}

\begin{algorithm}[htbp]
\caption{On-device parsing of end-to-end NMS output (with same-class dedup)}
\KwIn{model output $T$ of shape \texttt{[1,300,6]}; per-class thresholds $\tau(\cdot)$}
\KwOut{deduplicated detections $D$}
$D \leftarrow \emptyset$\;
\ForEach{row $r=(x_1,y_1,x_2,y_2,\mathit{conf},\mathit{cls}) \in T$}{
  \If{$\mathit{conf} \geq \tau(\mathit{cls})$ \textbf{and} box geometry is valid}{
    $c \leftarrow$ center of $r$\;
    \If{some kept $r' \in D$ has class $\mathit{cls}$ and contains $c$}{
      replace $r'$ with $r$ if $\mathrm{conf}(r) > \mathrm{conf}(r')$ \tcp*{in-graph NMS leaks identical rows; dedup required}
    }
    \Else{add $r$ to $D$}
  }
}
\Return $D$\;
\end{algorithm}

\section{Experimental Design}

Detection quality is measured with the standard object-detection metrics. With true
positives (TP), false positives (FP), and false negatives (FN), and predicted/ground-truth
boxes $A,B$:
\begin{equation}\mathrm{IoU}(A,B)=\frac{|A\cap B|}{|A\cup B|}\end{equation}
\begin{equation}\mathrm{Precision}=\frac{\mathrm{TP}}{\mathrm{TP}+\mathrm{FP}}\end{equation}
\begin{equation}\mathrm{Recall}=\frac{\mathrm{TP}}{\mathrm{TP}+\mathrm{FN}}\end{equation}
\begin{equation}\mathrm{AP}=\int_{0}^{1} p(r)\,\mathrm{d}r\end{equation}
\begin{equation}\mathrm{mAP}=\frac{1}{N}\sum_{i=1}^{N}\mathrm{AP}_i\end{equation}
where $p(r)$ is the precision--recall curve and AP its area; mAP@0.5 evaluates at
IoU $=0.5$ and mAP@0.5:0.95 averages AP over IoU thresholds $0.5,0.55,\dots,0.95$. We
also report the number of \emph{broken} classes (mAP@0.5 $<0.30$), since per-class
failure matters more than the average for a visitor-facing app.

\section{Results and Analysis}

Table~\ref{tab:iteration} summarizes the four iterations and Figure~\ref{fig:iteration}
plots the mAP@0.5 progression. The dominant lesson is that label cleaning ($+0.116$
mAP@0.5, no model change) and hand annotation ($+0.128$ mAP@0.5, $+0.147$ mAP@0.5:0.95)
broke ceilings that resolution scaling could not.

\begin{table}[htbp]
\centering\small
\setlength{\tabcolsep}{5pt}
\caption{Detector iteration comparison (validation).}
\label{tab:iteration}
\begin{tabular}{@{}lllllll@{}}
\toprule
\textbf{Iter.} & \textbf{Labels} & \textbf{imgsz} & \textbf{Head} & \textbf{mAP$_{50}$} & \textbf{mAP$_{50:95}$} & \textbf{Broken}\\
\midrule
v1     & Auto (YOLO-World)       & 320 & raw     & 0.751 & ---   & 10\\
v2     & Auto $+$ 4 rules        & 320 & raw     & 0.867 & 0.586 & 3\\
v2@640 & Auto $+$ rules          & 640 & E2E NMS & 0.867 & 0.777 & 3\\
\textbf{v3} & \textbf{Hand-annotated} & \textbf{640} & \textbf{E2E NMS} & \textbf{0.995} & \textbf{0.924} & \textbf{0}\\
\bottomrule
\end{tabular}
\end{table}

\begin{figure}[htbp]
\centering
\begin{tikzpicture}
\begin{axis}[ybar,symbolic x coords={v1,v2@320,v2@640,v3},xtick=data,
  ymin=0,ymax=1.08,ylabel={mAP@0.5},bar width=16pt,width=0.82\textwidth,height=6cm,
  nodes near coords,every node near coord/.append style={font=\scriptsize},
  enlarge x limits=0.18]
\addplot[fill=gemblue!55] coordinates {(v1,0.751)(v2@320,0.867)(v2@640,0.867)(v3,0.995)};
\end{axis}
\end{tikzpicture}
\caption{Detection mAP@0.5 across iterations: label cleaning (v1$\rightarrow$v2) and
hand annotation (v2@640$\rightarrow$v3) drive the gains; resolution scaling alone leaves
mAP@0.5 flat.}
\label{fig:iteration}
\end{figure}

Hand annotation rescued every previously broken class
(Table~\ref{tab:rescue}). Figure~\ref{fig:detexample} shows live detections from the
deployed app, including the Amenhotep~II / Thutmose~III sibling pair distinguished
simultaneously.

\begin{table}[htbp]
\centering
\caption{Per-class rescue from auto-labels (v2) to hand-labels (v3), mAP@0.5:0.95.}
\label{tab:rescue}
\begin{tabular}{@{}lcc@{}}
\toprule
\textbf{Class} & \textbf{v2} & \textbf{v3}\\
\midrule
Double Statue of Ramesses II        & 0.000 & 0.995\\
Overseer Amenemhat                  & 0.000 & 0.940\\
Thutmose III Statue                 & 0.000 & 0.941\\
Naktmin and Tiy                     & 0.269 & 0.977\\
Statue of King Akhenaten            & 0.281 & 0.955\\
Seated Statue of King Thutmose III  & 0.440 & 0.862\\
\bottomrule
\end{tabular}
\end{table}

\begin{figure}[htbp]
\centering
\begin{minipage}{0.40\textwidth}\centering
\includegraphics[width=\linewidth]{app_detection_1.jpg}\\[1mm](a)
\end{minipage}\hfill
\begin{minipage}{0.40\textwidth}\centering
\includegraphics[width=\linewidth]{app_detection_2.jpg}\\[1mm](b)
\end{minipage}
\caption{Live on-device detection in the deployed \TimeLens{} app: (a) two
near-identical New Kingdom statues distinguished at once---Amenhotep~II versus
Thutmose~III, the canonical fine-grained sibling case; (b) the Kneeling Statue of Queen
Hatshepsut. Each box appears only after clearing the on-device confidence gate.}
\label{fig:detexample}
\end{figure}

\section{Failure Modes and Caveats}

The v3 validation metrics \emph{overstate} real-world generalization, and we disclose
this plainly.

\textbf{Frame leakage.} The dataset is video frames split randomly per image. Adjacent
frames are near-duplicates, so a frame in training can have an almost-identical twin in
validation. The v3 mAP@0.5 saturates at $0.995$ within about three epochs of
warm-started training, which is characteristic of memorization rather than
generalization. A conservative, in-distribution \emph{held-out test} (using the
$320$ model) gives a more honest picture: precision $0.691$, recall $0.815$,
mAP@0.5 $0.751$, mAP@0.5:0.95 $0.650$.

\textbf{Domain shift.} Training frames come from stable, pre-opening videos; deployment
uses handheld phone cameras under public lighting with motion blur. Phone-camera
inference is materially worse than on the source videos.

\textbf{Auto-annotation ceiling.} Open-vocabulary detectors handle generic concepts but
fail on fine-grained, domain-specific objects; spatial rules recovered most failures,
but a residual fraction required manual labeling. Figure~\ref{fig:confusion} shows the
residual fine-grained confusion as an off-diagonal structure in the confusion matrix.

\begin{figure}[htbp]
\centering
\includegraphics[width=0.82\textwidth]{confusion_matrix.png}
\caption{Normalized 51-class confusion matrix (v3). The strong diagonal confirms
separation; faint off-diagonal cells mark residual sibling-class confusion among
visually similar statues.}
\label{fig:confusion}
\end{figure}

\section{Engineering Findings}

\begin{itemize}[leftmargin=1.5em]
  \item \textbf{Android frame rotation.} Back-camera sensors deliver frames rotated
        $90^{\circ}$ from the display; a rotation step is required before resizing or
        every prediction is computed on a sideways artifact.
  \item \textbf{End-to-end NMS format and duplicate leakage.} The exported output has shape
        \texttt{[1,300,6]}; each row is $(x_1,y_1,x_2,y_2,\mathit{conf},\mathit{cls})$ in
        corner format (not center-width-height), so misreading it stretches and shifts boxes.
        Empirically, the Ultralytics \texttt{nms=True} TFLite export bakes near-zero NMS
        thresholds into the graph, so a single dominant object can produce $100$--$250$
        identical copies of its top-survivor row in the $300$-slot buffer. The on-device
        parser therefore runs a per-class same-center dedup before drawing; without it,
        the overlay would draw the same box over itself many times.
  \item \textbf{Calibration from clean labels.} Clean labels produce sharp confidence
        distributions (real artifacts cluster high; off-target scenes fall well below),
        which lets the app raise its on-device threshold for out-of-distribution
        rejection (global $0.75$, with $0.90$ for the false-positive-prone
        \emph{Siptah Stela}).
  \item \textbf{Footprint.} The deployed asset is FP16 TFLite at \textbf{5.97\,MB},
        running in roughly $30$--$60$\,ms per frame on a mid-tier Android phone.
\end{itemize}

\section{Limitations}

The headline validation metrics are bounded by frame leakage, and a training-to-phone
domain gap remains. The honest, deployment-facing numbers are the held-out test figures
above. Mitigations are deferred to future work (Chapter~8): splitting by video rather
than by frame, capturing a phone-camera deployment test set, and stronger augmentation.

\section{Conclusion}

For fine-grained museum artifacts, \textbf{label quality is first-order}; architecture
and resolution are second-order. A compact YOLOv8n trained on clean, hand-verified
labels yields a deployable, real-time, on-device detector that feeds the rest of the
\TimeLens{} system.
